\documentclass{beamer}

\usepackage{beamerthemesplit}
\beamertemplatenavigationsymbolsempty
\setbeamertemplate{footline}[page number]{}

\usepackage{graphicx}
\usepackage{color}
\usepackage{listings}
\usepackage{verbatim}
\usepackage{setspace}
\usepackage{url}

\usepackage{fontspec}
\setmonofont{Latin Modern Mono}
\setsansfont{TeX Gyre Heros}

\lstset{breakatwhitespace=true,
language=C++,
basicstyle=\footnotesize\ttfamily,
keywordstyle=\color{blue}\ttfamily,
stringstyle=\color{red}\ttfamily,
commentstyle=\color{brown}\ttfamily,
morecomment=[l][\color{magenta}]{\#}keepspaces=true,
breaklines=true,
tabsize=3,
showstringspaces=false,
extendedchars=true,
frame=single}

\newcommand*{\vcenteredhbox}[1]{\begingroup
\setbox0=\hbox{#1}\parbox{\wd0}{\box0}\endgroup}


\title{Bio-Formats C++ Conversion}
\subtitle{Differences between C++ and Java}
\author{Roger Leigh}
\date{March 2014\\University of Dundee}

\begin{document}

\begin{frame}[plain]
  \titlepage
  \begin{center}
    \vcenteredhbox{\includegraphics[width=0.25\textwidth]{ome}} \hfill
    \vcenteredhbox{\includegraphics[width=0.2\textwidth]{dundee}}\hfill
    \vcenteredhbox{\includegraphics[width=0.25\textwidth]{wellcome}}
  \end{center}
\end{frame}

\section*{Overview}
\frame{
\frametitle{Overview}
\tableofcontents
}

\section{C++ and Java differences}
\subsection{C++ and Java types}

\begin{frame}[fragile]
  \frametitle{Java types}
  \begin{lstlisting}[language=Java]
int i;
double d;

Pixels pixels = new Pixels();

Pixels[] array = new Pixels[5];
\end{lstlisting}
  \begin{itemize}
    \pause
  \item Primitive types
    \pause
  \item No unsigned integer types
    \pause
  \item Classes, all derived from root \texttt{Object}
    \pause
  \item Objects are by reference only
    \pause
  \item Objects and arrays are always allocated with \texttt{new}
    \pause
  \item Destruction is non-deterministic
  \end{itemize}
\end{frame}

\begin{frame}[fragile]
  \frametitle{C++ types}
  \begin{lstlisting}
uint32_t i;
double d;

Pixels         pixels;
Pixels        *pixelsptr1 = new Pixels();
const Pixels  *pixelsptr2 = &pixels;
Pixels&        pixelsref(pixels);

Pixels array[5];
\end{lstlisting}
  \begin{itemize}
    \uncover<1->
  \item Primitive types, including unsigned, of standard sizes
    \uncover<1->
  \item Classes have no common root
    \uncover<1->
  \item All types may be instances, pointers or references
    \pause
  \item Pointers and references may refer to \texttt{const} type
    \pause
  \item Pointers may be \texttt{const}
    \pause
  \item Destruction is deterministic
  \end{itemize}
\end{frame}

\subsection{Reference handling}

\begin{frame}[fragile]
  \frametitle{C++ pointers: pitfalls of “dumb” pointers}
  \begin{lstlisting}
{
  Image *i = new Image(filename);

  i->read_plane(); // throws exception

  // Memory leaked here
  // Memory not freed when pointer i goes out of scope
}
\end{lstlisting}
  \begin{itemize}
    \item \texttt{new} and \texttt{delete} must always be paired
    \item Requires manual management of ownership and lifetime
    \item Bare pointers are not exception-safe
    \item Need to clean up for every exit point in a function
  \end{itemize}
\end{frame}

\begin{frame}[fragile]
  \frametitle{C++ pointers: \texttt{std::shared\_ptr} as a “smart” pointer}
  \begin{lstlisting}
{ // Image *i = new Image(filename);
  std::shared_ptr<Image> i(std::make_shared<Image>(filename));

  i->read_plane(); // throws exception

  // Memory freed when i's destructor is
  // run at exit of block scope
}

\end{lstlisting}
  \begin{itemize}
    \item \texttt{shared\_ptr} object lifetime manages the resource
    \item May be used as class members; lifetime is tied to class instance
    \item Clean up for all exit points is automatic and safe
    \item Allows ownership transfer and sharing
    \item Allows reference without ownership using \texttt{weak\_ptr}
  \end{itemize}
\end{frame}

\begin{frame}[fragile]
  \frametitle{RAII: Resource Acquisition Is Initialisation}
    \begin{itemize}
    \item Class is a proxy for a resource
    \item Resource is acquired when object is intialised
    \item Resource is released when object is destroyed
    \item Manage any resource (memory, files, locks, mutexes)
    \item C++ language and runtime guarantees make resource management
      deterministic and reliable
    \item Safe for use in any scope
    \item Exception safe
    \item Used throughout modern C++ libraries and applications
    \end{itemize}
\end{frame}

\begin{frame}[fragile]
  \frametitle{C++ reference variants}
  \begin{lstlisting}[basicstyle=\tiny\ttfamily]
//          Non-constant                              Constant
// -----------------------------       ---------------------------------------
// Pointer
                       Image *i;                               const Image *i;
                Image * const i;                        const Image * const i;
§\pause§
// Reference
                       Image& i;                               const Image& i;
§\pause§
// Shared pointer
       std::shared_ptr<Image> i;               std::shared_ptr<const Image> i;
 const std::shared_ptr<Image> i;         const std::shared_ptr<const Image> i;
§\pause§
// Shared pointer reference
      std::shared_ptr<Image>& i;              std::shared_ptr<const Image>& i;
const std::shared_ptr<Image>& i;        const std::shared_ptr<const Image>& i;
§\pause§
// Weak pointer
         std::weak_ptr<Image> i;                 std::weak_ptr<const Image> i;
   const std::weak_ptr<Image> i;           const std::weak_ptr<const Image> i;
§\pause§
// Weak pointer reference
        std::weak_ptr<Image>& i;                std::weak_ptr<const Image>& i;
  const std::weak_ptr<Image>& i;          const std::weak_ptr<const Image>& i;
\end{lstlisting}
\pause
\begin{center}
  Java has one reference type.  Here, we have \textbf{22}…
\end{center}
\end{frame}

\begin{frame}[fragile]
  \frametitle{C++ reference usage rationalised}
\begin{block}{Class members}
  \begin{lstlisting}[basicstyle=\scriptsize\ttfamily]
Image i;                     // Concrete instance§\pause§
std::shared_ptr<Image> i;    // Reference§\pause§
std::weak_ptr<Image> i;      // Weak reference
\end{lstlisting}
\pause
\end{block}
\begin{block}{Arguments}
  \begin{lstlisting}[basicstyle=\scriptsize\ttfamily]
void read_plane(const Image& image);§\pause§
void read_plane(const std::shared_ptr<Image>& image);
\end{lstlisting}
\pause
\end{block}
\begin{block}{Return}
  \begin{lstlisting}[basicstyle=\scriptsize\ttfamily]
                  Image get_image(); // ownership transferred§\pause§
                 Image& get_image(); // ownership retained§\pause§
 std::shared_ptr<Image> get_image(); // ownership shared/trans§\pause§
std::shared_ptr<Image>& get_image(); // ownership shared
\end{lstlisting}
\end{block}
\end{frame}

\begin{frame}[fragile]
  \frametitle{C++ reference usage rationalised}
  \begin{itemize}
  \item<1-> Safety: References can not be \texttt{null}.
  \item<2-> Storing polymorphic types requires use of a
    \texttt{shared\_ptr}.
  \item<3-> Safety: To avoid cyclic dependencies, use \texttt{weak\_ptr}.
  \item<4-> Safety: To allow object destruction while maintaining a safe
    reference, use \texttt{weak\_ptr}.
  \item<5-> \texttt{weak\_ptr} is not directly usable.
  \item<5-> \texttt{weak\_ptr} is convertible back to \texttt{shared\_ptr}
    for use \emph{if the object is still in existence}.
  \item<6-> C++11 \emph{move semantics} (\texttt{\&\&}) improve
    performance of ownership transfer
  \end{itemize}
\end{frame}

\subsection{Additional API concerns}

\begin{frame}[fragile]
  \frametitle{Safe array passing: \texttt{std::array} / \texttt{boost::array}}
  \begin{lstlisting}
class Image
{
  // Unsafe; size unknown
  uint8_t[] getLUT();
       void setLUT(uint8_t[]& lut);

  // Safe; size defined
  typedef std::array<uint8_t, 256> LUT;
  const LUT& getLUT() const;
        void setLUT(const LUT&);
};
\end{lstlisting}
  \begin{itemize}
  \item<1-> C (C++) arrays essentially decay to plain pointers
  \item<2-> Plain pointers have no associated size information
  \item<3-> \texttt{std::array} is an array-like object
  \item<4-> \texttt{std::array} size defined in the template
  \item<5-> \texttt{std::array} can be passed like any object
  \end{itemize}
\end{frame}

\begin{frame}[fragile]
  \frametitle{Containers storing different types}
\begin{block}{Types with a common base}
  \begin{lstlisting}
    std::vector<std::shared_ptr<Base> > v;
    v.push_back(std::make_shared<Derived>());
\end{lstlisting}
  \begin{itemize}
  \item Store any type derived from \texttt{Base}
  \end{itemize}
\end{block}
\begin{block}{Java containers can be problematic}
  \begin{itemize}
  \item Java can store root \texttt{Object} in containers
  \item Java can pass and return root \texttt{Object} in methods.
  \item This isn't possible in C++: there is no root object.
  \item An alternative approach is needed.
  \end{itemize}
\end{block}
\end{frame}

\begin{frame}[fragile]
  \frametitle{Containers storing different types}
\begin{block}{Types with no common base: \texttt{boost::variant}}
  \begin{lstlisting}
    typedef boost::variant<int, std::string> variants;
    std::vector<variants> v;
    v.push_back(43);
    v.push_back("ATTO 647N")
\end{lstlisting}
  \begin{itemize}
  \item Store a set of discriminated types
  \item Used to store original metadata
  \end{itemize}
\end{block}
\begin{block}{Arbitrary types: \texttt{boost::any}}
  \begin{lstlisting}
    std::vector<boost::any> v;
    v.push_back(Anything);
\end{lstlisting}
  \begin{itemize}
  \item Flexible, but need to cast to each type used to extract.
  \end{itemize}
\end{block}
\end{frame}

\begin{frame}[fragile]
  \frametitle{Interfaces}
\begin{block}{Java: Single inheritance, plus interfaces}
\end{block}
\begin{block}{C++: Multiple inheritance}
  \begin{itemize}
  \item Interfaces are classes with:
    \begin{itemize}
      \item No instance variables
      \item Pure virtual methods
      \item \texttt{protected} default constructor
      \item \texttt{public virtual} destructor
      \item Deleted copy constructor and assignment operator
    \end{itemize}
  \item Classes implementing interfaces:
    \begin{itemize}
      \item Use \texttt{public} inheritance for parent class
      \item Use \texttt{virtual public} inheritance for implemented interfaces
      \item \texttt{virtual} destructor
    \end{itemize}
  \end{itemize}
\end{block}
\end{frame}

\section*{Acknowledgements}

\frame{
  \frametitle{Acknowledgements}
  \parbox[t]{0.45\textwidth}{
    \begin{itemize}
    \item OME Team, Dundee
      \begin{itemize}
      \item Jason Swedlow
      \item Jean-Marie Burel
      \item Mark Carroll
      \item Andrew Patterson
      \item …and the rest of the team
      \end{itemize}
    \end{itemize}
  }
  \parbox[t]{0.45\textwidth}{
    \begin{itemize}
    \item Micron, Oxford
      \begin{itemize}
      \item Douglas Russell
      \end{itemize}
    \item Glencoe Software
      \begin{itemize}
      \item Melissa Linkert
      \item Josh Moore
      \end{itemize}
    \end{itemize}
  }

  \begin{center}
    \vcenteredhbox{\includegraphics[width=0.25\textwidth]{ome}} \hfill
    \vcenteredhbox{\includegraphics[width=0.2\textwidth]{dundee}}\hfill
    \vcenteredhbox{\includegraphics[width=0.5\textwidth]{wellcome}}
  \end{center}
}

\end{document}
